\documentclass[handout_subfiles_root.tex]{subfiles}
\usepackage[rm2]{lecturenote}

\begin{document}
\HandoutTitle{1}{Signals: Analog, Discrete-Time, and Digital}{August 25, 2026}
\maketitle

\section{What is a Digital Signal? (1D Case)}

We distinguish three flavors of one-dimensional signals, based on how the
independent variable (time) and the dependent variable (amplitude) are
represented.

\subsection{Continuous-Time (``Analog'') Signals}

A continuous-time signal $x_a(t)$ is a real- or complex-valued function
defined for all real-valued $t$.

\begin{center}
\begin{tikzpicture}[scale=1]
\draw[-Latex] (-0.3,0) -- (6.6,0) node[right]{$t$};
\draw[thick, ee483blue, domain=0:6.3, samples=150]
  plot (\x, {0.9*sin(deg(\x*3.4))*exp(-0.12*\x) + 0.05*\x - 0.05});
\end{tikzpicture}
\end{center}

\noindent Examples: speech signals, EEG, images formed with focused light.

\subsection{Discrete-Time Signals}

A discrete-time signal $x[n]$ is defined only for integer-valued $n$.

\begin{center}
\begin{tikzpicture}[scale=1]
\draw[-Latex] (-0.5,0) -- (5.5,0) node[right]{$n$};
\foreach \n/\val/\lbl in {0/1.2/{\pi/2},1/1.7/{2},2/-1.0/{-1},3/1.4/{},4/1.0/{}}
  {\draw[thick] (\n,0) -- (\n,\val); \fill (\n,\val) circle (2.2pt);}
\node[above] at (0,1.25) {\scriptsize $\pi/2$};
\node[above right] at (1,1.75) {\scriptsize $2$};
\node[below] at (2,-1.05) {\scriptsize $-1$};
\foreach \n in {0,1,2,3,4} \node[below] at (\n,-0.15) {\n};
\end{tikzpicture}
\end{center}

Discrete-time signals arise in (at least) two ways:
\begin{itemize}
  \item \textbf{Sampled continuous signals.} $x[n] = x_a(nT)$, e.g.\ a
  sampled audio signal.
  \item \textbf{Inherently discrete signals.} Quantities that only update at
  discrete instants, e.g.\ the daily maximum temperature, or the number of
  students late to class on days when class meets.
\end{itemize}

\subsection{Digital Signals}

\begin{defbox}[title={Definition (Digital Signal)}]
A \textbf{digital signal} is sampled in time \emph{and} in amplitude: each
value $x[n]$ is stored with a finite number of bits, so the numbers are
rounded. This is called \textbf{finite precision} or \textbf{quantization}.
\end{defbox}

\section{From the Real World to Numbers}

\subsection{The Real ADC}

\begin{center}
\begin{tikzpicture}[node distance=14mm, >=Latex]
\node (in) {$x_a(t)$};
\node[draw, minimum width=2.4cm, minimum height=1cm, right=of in] (adc) {ADC};
\node[right=of adc] (out) {$q[n]$};
\draw[->] (in) -- (adc);
\draw[->] (adc) -- (out);
\node[below=2pt of adc, font=\scriptsize] {analog-to-digital converter};
\node[above=2pt of out, font=\scriptsize, align=left] {digitized signal\\ (quantized in amplitude)};
\end{tikzpicture}
\end{center}

\subsection{A Useful Model}

We conceptually split the ADC into a time-sampling stage followed by an
amplitude-quantization stage:

\begin{center}
\begin{tikzpicture}[node distance=12mm, >=Latex]
\node (in) {$x_a(t)$};
\node[draw, minimum width=1.9cm, minimum height=1cm, right=of in, align=center] (samp) {time\\ sample};
\node[right=of samp] (xn) {$x[n]$};
\node[draw, minimum width=1.9cm, minimum height=1cm, right=of xn] (quant) {quantize};
\node[right=of quant] (qn) {$q[n]$};
\draw[->] (in) -- (samp);
\draw[->] (samp) -- (xn);
\draw[->] (xn) -- (quant);
\draw[->] (quant) -- (qn);
\node[below=2pt of samp, font=\scriptsize] {discrete-time};
\end{tikzpicture}
\end{center}

\subsection{An Even More Useful Model}

Quantization is modeled as \emph{additive quantization noise} $e[n]$ injected
after ideal time-sampling:

\begin{center}
\begin{tikzpicture}[node distance=12mm, >=Latex]
\node (in) {$x_a(t)$};
\node[draw, minimum width=1.9cm, minimum height=1cm, right=of in, align=center] (samp) {time\\ sampling};
\node[right=of samp] (xn) {$x[n]$};
\node[draw, circle, right=of xn] (sum) {$+$};
\node[right=of sum] (qn) {$q[n]$};
\node[above=8mm of sum] (e) {quantization noise $e[n]$};
\draw[->] (in) -- (samp);
\draw[->] (samp) -- (xn);
\draw[->] (xn) -- (sum);
\draw[->] (e) -- (sum);
\draw[->] (sum) -- (qn);
\end{tikzpicture}
\end{center}

\begin{notebox}
Discrete-time systems \emph{can} be linear. \textbf{Digital systems cannot be
linear}; quantization is an inherently nonlinear operation.
\end{notebox}

\section{DSP Topics for the Course}

\begin{itemize}
  \item \textbf{Sampling and reconstruction} (ADC and DAC)
  \begin{itemize}
    \item Nyquist rate
  \end{itemize}
  \item \textbf{Discrete-time linear systems}
  $$x_a(t) \to \boxed{\text{linear system}} \to y_a(t), \qquad
    x[n] \to \boxed{\text{linear system}} \to y[n]$$
  \item \textbf{Signal representation and coding}
  \begin{itemize}
    \item Discrete Time Fourier Transform (DTFT)
    \item Discrete Fourier Transform (DFT)
    \item Discrete cosine transform and wavelets (WebP, JPEG, MPEG, MP3, etc.)
  \end{itemize}
  \item \textbf{Implementation}
  \begin{itemize}
    \item Design of digital filters
    \item Hardware (CPU, DSP, GPU, FPGA)
    \item Computational cost (algorithmic) and quantization / overflow
    (configuration details)
  \end{itemize}
\end{itemize}

\section{Why DSP Instead of Analog Signal Processing?}

\begin{itemize}
  \item Cost
  \item Storage of data
  \item Availability
  \item Easily reconfigured
  \item Implementation time
  \item Reliability
  \item Flexibility
\end{itemize}

\section*{Optional Readings}
\begin{itemize}
  \item Sanjit K. Mitra, \textit{Digital Signal Processing: A Computer-Based Approach}, 4th ed., Ch.\ 1, 2.1, 2.2, 2.4
  \item Monson H. Hayes, \textit{Schaum's Outline of Digital Signal Processing}, 2nd ed., Ch.\ 1.1--1.3
\end{itemize}

\HandoutAppendixListStart
  \HandoutAppendixListItem{app:scan}{A.\ Handwritten Lecture Note Scan (August 25, 2026)}
\HandoutAppendixListEnd

\HandoutAppendixSection{app:scan}{Handwritten Lecture Note Scan}
\HandoutIncludeHandoutScan{lecture1_0825.pdf}

\end{document}
